% DRAFT, 8 Nov 2026 (revised after an adversarial AI review the same day, and on 9 Nov after
% the independent referee report 2026-11-09-rees-blackbox-corollary-referee.md, findings
% 1-6, 8, 9 applied; finding 7, core self-containment, is an owner decision).
% AI-written, AI-checked only, pending human review. NOT yet in the paper: it waits
% for the single paper-update pass (owner, 7 Nov: "wait until we have all").
%
% Prompted by an outside question (Nov 2026) about how strong the explicit-chart
% assumption is.
%
% Where it goes, and what else must change in the same pass:
%  * (A) after Definition~\ref{def:input-contract} in core/main-core.tex, and in main.tex.
%  * (B) after Theorem~\ref{thm:main} and its remarks in core/main-core.tex. In main.tex,
%    state it as a corollary of Proposition~\ref{prop:black-box}: it is the e = 1
%    "Alternatively" route in the proof of Proposition~\ref{prop:black-box}(ii)
%    (appendix-torsion-blackbox.tex, near line 825), stated with Schedule~1 alone.
%  * core/main-core.tex l.184 ("we do not claim it here") and the outlook paragraph
%    "Black-box nilpotent groups" (l.697, "not claimed here") must be revised to match:
%    exponent-p p-groups become claimed (conditional on Theorem~\ref{thm:main}, like
%    everything else), the nilpotent composition and exponent p^e stay unclaimed.
%  * The core does not define \UT (main.tex:35 does), has no bib entry
%    CicaloDeGraafVaughanLee2012, and cannot \ref the full version's labels; full-version
%    results are cited in words below ("Lemma ... of \fullv").
%  * The core's policy is that no proof in it depends on the full version. The proof
%    below uses two lemmas proved there (tb-evaluation, tb-kernel); both are standard
%    (inverse Hausdorff formulas; abelian Fourier sampling), cited from the literature
%    as well. NEEDS HUMAN VERIFICATION: the Khukhro chapter for the inverse formulas.
%
% Finite control: theorems/rees_blackbox_end_to_end.py,
% research/audits/reviews/2026-11-08-rees-blackbox-e2e/, tests/test_rees_blackbox_end_to_end.py.

%----------------------------------------------------------------------
% (A) After Definition def:input-contract
%----------------------------------------------------------------------

\paragraph{How strong is the explicit chart?}
Definition~\ref{def:input-contract} asks for the structure constants of the Lazard Lie algebra. Three common presentations reduce to it.
\begin{enumerate}[label=(\alph*),leftmargin=2em]
  \item \emph{Unitriangular matrix groups.} Let $G \le \UT_n(\Fp)$ be given by generator matrices $g_1,\dots,g_s$ and an oracle $f$ on matrices, and suppose that the non-unital associative algebra $\mathcal{A}$ generated by the $g_i - I$ satisfies $\mathcal{A}^p = 0$. This holds whenever $n \le p$, since products of $n$ strictly upper triangular matrices vanish, and in general it can be checked by linear algebra in $\poly(n,s,\log p)$. Then $G \subseteq I + \mathcal{A}$, $G$ has class below $p$, and $G$ has exponent dividing $p$, since $(I+N)^p = I + N^p = I$. The exponential and the logarithm $\log(I+N) = \sum_{k=1}^{p-1}(-1)^{k+1}N^k/k$, truncated below degree $p$, are inverse bijections between $I+\mathcal{A}$ and $\mathcal{A}$, and their denominators are prime to $p$. In the free associative algebra over $\mathbb{Z}_{(p)}$ modulo degree $\ge p$, the truncated Baker--Campbell--Hausdorff series also has denominators prime to $p$, so $\exp X \exp Y = \exp \mathrm{BCH}(X,Y)$ on $\mathcal{A}$; by uniqueness of the Lazard ring structure~\cite[Chapters~9--10]{Khukhro1998}, $\log$ restricted to $G$ is the Lazard logarithm. The Lie subalgebra $\mathcal{L}'$ generated by the $\log g_i$ equals $\log G$: it lies in $\log G$, which is a Lie subalgebra, and $\exp\mathcal{L}'$ is a subgroup containing every $g_i$. Bracket closure with row reduction gives a basis of $\mathcal{L}'$ and its structure constants in time $\poly(n,s,\log p)$, and $y \mapsto f(\exp y)$ is the oracle of Definition~\ref{def:input-contract}, at one query of $f$ per query. Theorem~\ref{thm:main} then applies when the class is fixed. This covers the groups $\Fp \ltimes \Fp^m$ of~\cite{BaconChildsVanDam2005,IvanyosSantha2015} whenever $p$ exceeds the class $c$: the generator of $\Fp$ acts by a unipotent matrix, which a flag of kernels of its powers minus $I$ conjugates to upper unitriangular form in polynomial time, and in the resulting embedding into $\UT_{m+1}(\Fp)$ products of $c+1$ of the matrices $g_i - I$ vanish, so $\mathcal{A}^{c+1} = 0$.
  % AI-checked exactly for p = 3, m = 4, class 2 (order 243): research/audits/reviews/2026-11-08-rees-blackbox-e2e/semidirect_check.py.
  \item \emph{Black-box groups.} Corollary~\ref{cor:black-box} below acquires the chart from a black-box $p$-group of class below $p$ and exponent $p$ with unique encoding and a coherent multiplication oracle, with no query to $f$.
  \item \emph{Power-commutator presentations.} Normal forms of a consistent power-commutator presentation are unique encodings. So for a group of exponent $p$ and class below $p$, whenever normal forms can be multiplied in polynomial time, the presentation gives a black box to which Corollary~\ref{cor:black-box} applies. We do not analyse the cost of collection. For a $p$-group of class below $p$ given by a consistent power-commutator presentation, Cical\`o, de Graaf and Vaughan-Lee compute the Lie ring structure on the polycyclic generators by evaluating truncated inverse Baker--Campbell--Hausdorff formulae in the group with the collection algorithm~\cite[\S7]{CicaloDeGraafVaughanLee2012}, and construct the group of a given Lie ring by the Baker--Campbell--Hausdorff formula. Their method is practical and is not given a complexity bound; the chart acquisition in Corollary~\ref{cor:black-box} evaluates the same kind of formulae, but in a black box.
  % CicaloDeGraafVaughanLee2012: AI-read 9 Nov (2026-11-09-rees-blackbox-citation-check.md); doi 10.1016/j.jalgebra.2011.11.031; add its bib entry.
\end{enumerate}

%----------------------------------------------------------------------
% (B) After Theorem thm:main
%----------------------------------------------------------------------

\begin{corollary}[Black-box groups of exponent $p$; AI-checked, pending human review]
\label{cor:black-box}
Fix $c \ge 1$ and assume Theorem~\ref{thm:is-import}. Let $G = \langle g_1,\dots,g_s\rangle$ be a black-box group with unique $\ell$-bit encodings and a coherent multiplication oracle $|x,y,z\rangle \mapsto |x,y,z\oplus xy\rangle$, promised to be a $p$-group of class at most $c < p$ and exponent $p$, with $p$ given. Let $f$ be a coherent hiding function of label width $w$ for an arbitrary $H \le G$, and let $\eps \in (0,1)$; put $L = \lfloor \ell/\log_2 p\rfloor$. A quantum algorithm returns generators of $H$ with probability at least $1-\eps$. It makes no query to $f$ until it has computed an explicit chart. It always queries $f$ at most $c(L+\tau)\prod_{d=2}^{c}B_d(L)$ times, in the notation of Theorem~\ref{thm:main}, a bound polynomial in $\ell$ and independent of $s$; on the event, of probability at least $1-\eps/4$, that the kernel step returns $\ker\phi$ exactly, it queries $f$ at most $Q_{\mathrm q}$ times, the count of Theorem~\ref{thm:main} at error $\eps/2$ for $m = \log_p|G|$ and $g_d = \dim\gamma_d(\log G)$. Its other cost is $\poly_c(s,\ell,w,\log p,\log(1/\eps))$ gates and black-box multiplications.
\end{corollary}

\begin{proof}
Put $F = F_{s,c}$, the free nilpotent Lie algebra of class $c$ over $\Fp$ on $x_1,\dots,x_s$, of dimension $d = \sum_{k=1}^c \frac1k\sum_{j\mid k}\mu(j)s^{k/j} \le \sum_{k \le c}s^k$, with its Lyndon basis of brackets. Since $G$ has exponent $p$ and class below $p$, $\log G$ is a Lie algebra over $\Fp$ of class at most $c$, so there is a unique homomorphism $\phi : F \to \log G$ with $x_i \mapsto \log g_i$. Its image is a Lie subalgebra containing every $\log g_i$, so by the subgroup form of the Lazard correspondence $\phi$ is onto, and $K = \ker\phi$ is an ideal with $F/K \cong \log G$. Write $\Phi = \exp\circ\phi$.

\emph{Evaluating $\Phi$.} The Lazard sum and bracket of $\log G$ are computed from group elements by fixed straight-line programs of products and powers $z \mapsto z^{r}$, $r \in \mathbb{Z}[1/c!]$, depending only on $c$ (the inverse Baker--Campbell--Hausdorff formulas, \cite[Chapters~9--10]{Khukhro1998} and \cite[\S2]{CicaloDeGraafVaughanLee2012}; explicit certified programs and their evaluation in $G$, with $z^{a/b} = z^{ab^{-1} \bmod p}$, are Lemmas~tb-words and~tb-evaluation of~\fullv). Since $z^{-1} = z^{p-1}$, only the multiplication oracle is used. So the images $\Phi(w_j)$ of the Lyndon basis are bracket programs in the $g_i$, and $\Phi(u)$ is the Lazard linear combination $\sum_j u_j\Phi(w_j)$, at $O_c(d\log p)$ multiplications. $\Phi$ is computed coherently and its workspace cleaned up by uncomputation. We assume, as for $U_f$ in the core, that the multiplication oracle acts as some unitary on strings that are not encodings; the identity's encoding is $g_1^p$, and multiplication by a classical constant $h$ is controlled by selecting $\mathrm{enc}(h)$ or $\mathrm{enc}(1)$ into the second register.

\emph{Kernel.} Since $\phi$ is additive and $\exp$ is a bijection, $\Phi(u) = \Phi(v)$ exactly when $u - v \in K$. So $u \mapsto \mathrm{enc}(\Phi(u))$ hides the subgroup $K$ of $(F,+) \cong \Fp^d$, and abelian Fourier sampling, with no query to $f$, finds $K$ from $t = \min(d,\lfloor\ell/\log_2 p\rfloor) + \lceil\log_p(8/\eps)\rceil$ samples with failure probability at most $\eps/4$, including the approximation of the Fourier transforms (Lemma~tb-kernel of~\fullv, at $\delta = \eps/4$). Two public checks follow: that $\Phi(v) = 1$ for a basis of $K'$, which fails whenever $K' \not\subseteq K$, and that $K'$ is an ideal of $F$. If a check fails, or if $d - \dim K' > L$, the algorithm aborts. With exact Fourier transforms every sample lies in $K^\perp$, so $K' \supseteq K$ and the first check detects every sampling failure. A deviation $K' \subsetneq K$ can come only from the approximation error and is not detected; it does not affect correctness: if $K'$ is an ideal, $F/K'$ is a Lie algebra of class at most $c$ over $\Fp$ and $E' = \Phi \circ \iota'$ below is a surjective homomorphism onto $G$, so the core still returns a basis whose images generate $H$, at the query count for dimension $d - \dim K' \le L$.

\emph{Chart.} Take the row-reduced basis of $K$ and let $\iota : \Fp^m \to F$ put coordinates on the non-pivot positions, $m = d - \dim K = \log_p|G|$. Projecting brackets of $F$ along $K$ gives the structure constants of $F/K \cong \log G$ in these coordinates, by $\poly(d)$ field operations, and $E(y) = \Phi(\iota y)$ is a group isomorphism from $(\Fp^m,\ast)$ onto $G$, computed from the $m$ precomputed images of the coordinate vectors at $O_c(m\log p)$ multiplications.

\emph{Core.} The chart and $f \circ E$ form an instance of Definition~\ref{def:input-contract}, and $f \circ E$ hides $E^{-1}(H)$. Theorem~\ref{thm:main} at error $\eps/2$ returns a basis of $\log E^{-1}(H)$, with $Q_{\mathrm q}$ queries of $f\circ E$, each one query of $f$ and two coherent evaluations of $E$. The images of the basis under $E$ generate $H$. A union bound over the kernel step ($\eps/4$) and the core ($\eps/2$) gives the error. The costs are polynomial because $d \le c\,s^c$ and $m\log_2 p \le \ell$; they include the bracket table of $F$ in its Lyndon basis, the ideal check ($\dim K' \cdot d$ brackets) and the push-forward of the output basis, each $\poly_c(s,\log p)$. The worst-case query bound holds because every branch that reaches the core has dimension at most $L$ and $g'_d \le L$, and $B_d$ is monotone.
\end{proof}

\begin{remark}[Scope of Corollary~\ref{cor:black-box}]
\label{rem:black-box-scope}
The corollary is a field form of the $e = 1$ chart route in the proof of Proposition~black-box(ii) of~\fullv, with a row-reduced complement in place of the chart lemma used there, stated with Schedule~1 alone. Like Theorem~\ref{thm:main}, it is conditional only on Theorem~\ref{thm:is-import}; its proof uses Lemmas~tb-words, tb-evaluation and tb-kernel of~\fullv. Given class below $p$, the exponent-$p$ promise can be tested, since $G$ has exponent $p$ exactly when every $g_i^p = 1$ ($O(s\log p)$ multiplications). Four restrictions remain. (i) Exponent $p$: for $p$-groups of class below $p$ and exponent $p^e > p$,~\fullv{} gives a route through a ring-to-field recursion that is conditional and is not claimed here. (ii) Unique encoding: the corollary, like~\cite{ImranIvanyos2023}, assumes unique encodings. (iii) $p$-groups: for black-box nilpotent groups of class at most $c$, the composition in~\fullv{} (its Theorem~bb-envelope) needs, with Schedule~1 alone and exponent-$p$ Sylow subgroups, beyond Theorem~\ref{thm:main} only a Sylow split by order finding and factoring and the exact algorithm of~\cite[Theorem~1]{ImranIvanyos2023} for the primes $p \le c$; this paper does not treat those steps, and the composition is not claimed here. (iv) Oracle input only: the coset-copy input of Theorem~\ref{thm:main} is not covered, since $E^{-1}$ is never computed.

A finite simulation (\file{theorems/rees_blackbox_end_to_end.py}; Appendix~\ref{app:artifacts}) runs the chart step of the corollary from opaque handles on three small unitriangular matrix groups and on the free two-generator class-three group of exponent $5$ given by hidden coordinates, and then the scaled exact simulator of the core on the acquired chart (with the per-grade arities of~\fullv, not Schedule~1 alone). The Fourier sampling law is computed from the fibres of $\Phi$ in floating point, and the homomorphism check on the chart is exhaustive or sampled as recorded. In every run the subgroup generated by the returned handles equals $H$. It is a finite check and proves nothing for other groups.
\end{remark}
